% Propietario preservado y actualizado por los cierres sucesores.
\section{Autoinercia como conexi\'on end\'ogena del estado conjunto}

Sea \(\nabla^{\rm tot}\) una conexi\'on sobre el espacio de estados
conjuntos. El marco local del observador se obtiene por restricci\'on:
\begin{equation}
 \mathfrak F_{\mathcal O}(\Gamma)
 =\operatorname{Res}_{\mathcal O}
 (\Gamma,\nabla^{\rm tot}).
 \label{eq:rm-obs-induced-frame}
\end{equation}
Una trayectoria completa \(\Gamma(t)\) es autoinercial cuando
\begin{equation}
 \nabla^{\rm tot}_{\dot\Gamma}\dot\Gamma=0.
 \label{eq:rm-obs-total-geodesic}
\end{equation}
Sea \(\pi_S\) la proyecci\'on al subsistema observado y
\(\gamma_S=\pi_S\circ\Gamma\). Definimos el defecto de proyecci\'on
\begin{equation}
 \mathcal K_{\pi_S}(\dot\Gamma,\dot\Gamma)
 :=\nabla^S_{\dot\gamma_S}\dot\gamma_S
 -\dd\pi_S\!\left(
 \nabla^{\rm tot}_{\dot\Gamma}\dot\Gamma
 \right).
 \label{eq:rm-obs-projection-defect}
\end{equation}
Sobre una trayectoria autoinercial queda
\begin{equation}
 \boxed{
 a_S=\nabla^S_{\dot\gamma_S}\dot\gamma_S
 =\mathcal K_{\pi_S}(\dot\Gamma,\dot\Gamma).}
 \label{eq:rm-obs-projected-acceleration}
\end{equation}
La aceleraci\'on del subsistema mide el acoplamiento y los grados de libertad
omitidos por la proyecci\'on. La totalidad posee una conexi\'on end\'ogena;
sus sombras pueden presentar aceleraci\'on, rotaci\'on y fuerza efectiva.

\begin{definition}[Estado autoinercial]
Un estado conjunto es autoinercial respecto de una conexi\'on y una familia
de observadores cuando su evoluci\'on satisface
\cref{eq:rm-obs-total-geodesic} y cada marco local se obtiene mediante
\cref{eq:rm-obs-induced-frame}. Las fuerzas aparentes de cada subsistema se
registran por los defectos \(\mathcal K_{\pi_S}\).
\end{definition}

La definici\'on integra observador y observable dentro de una sola din\'amica.
El cambio de marco corresponde a otra restricci\'on del mismo estado y de la
misma conexi\'on. La distinci\'on entre transporte af\'in, rotaci\'on de la
t\'etrada, tors\'ion y curvatura permanece disponible en el codominio
geom\'etrico.

\section{Criterio machiano de factorizaci\'on global}

Sea \(b:\mathcal X_\infty\to B_\infty\) la traza global de frontera y sea
\(I_v:\mathcal X_\infty\to V_v\) una respuesta inercial local. La
respuesta es machiana cuando existe
\begin{equation}
 \overline I_v:\operatorname{im}b\longrightarrow V_v,
 \qquad
 I_v=\overline I_v\circ b.
 \label{eq:rm-obs-mach-factorization}
\end{equation}

\begin{theorem}[Criterio de respuesta machiana]
\label{thm:rm-obs-mach-criterion}
Existe una \`unica aplicaci\'on \(\overline I_v\) que satisface
\cref{eq:rm-obs-mach-factorization} si y s\'olo si \(I_v\) es constante
sobre cada fibra de \(b\):
\begin{equation}
 b(x)=b(x')\quad\Longrightarrow\quad I_v(x)=I_v(x').
 \label{eq:rm-obs-mach-fiber-criterion}
\end{equation}
Si \(b\) es sobreyectiva, \(\overline I_v\) queda definida en todo
\(B_\infty\).
\end{theorem}

\begin{proof}
La factorizaci\'on implica constancia en las fibras. Rec\'iprocamente, para
\(\beta\in\operatorname{im}b\), definimos
\(\overline I_v(\beta)=I_v(x)\) con \(x\in b^{-1}(\beta)\). La constancia
garantiza independencia del representante y la ecuaci\'on de factorizaci\'on
garantiza unicidad.
\end{proof}

El criterio traduce el principio de Mach en una propiedad verificable: la
respuesta local es una lectura de la clase global de cierre. Para demostrar
dependencia genuinamente global se exhiben adem\'as dos estados con igual
restricci\'on local y distinta traza de frontera cuya respuesta difiere.
Esta condici\'on separa una factorizaci\'on global efectiva de una respuesta
determinada exclusivamente por la vecindad local.

La escala gravitatoria \(Q\) del reloj CT108 proporciona una realizaci\'on
adimensional de esta lectura:
\begin{equation}
 \alpha_G(E_1,E_2)=\frac{E_1E_2}{Q^2}.
 \label{eq:rm-obs-global-clock-coupling}
\end{equation}
El reloj global fija \(Q\), cada secci\'on local publica \(E/Q\), y la
geometr\'ia responde mediante la compliancia \(G/c^4\). La ecuaci\'on
proporciona una carta machiana de la respuesta gravitatoria; el
\cref{thm:rm-obs-mach-criterion} determina el requisito de factorizaci\'on
que debe satisfacer su realizaci\'on sobre el estado completo.

\section{Holonom\'ia y jerarqu\'ia de retornos}

Sea \(\gamma\) un lazo basado en \(x_0\), y sea
\begin{equation}
 \mathscr H_\gamma=\Hol_\nabla(\gamma)
 \label{eq:rm-obs-holonomy}
\end{equation}
la holonom\'ia en la fibra \(E_{x_0}\). Un dato transportado retorna como
vector exactamente cuando pertenece a
\begin{equation}
 \operatorname{Fix}(\mathscr H_\gamma)
 =\{\psi:\mathscr H_\gamma\psi=\psi\}.
 \label{eq:rm-obs-fix-holonomy}
\end{equation}
La interpretaci\'on autoinercial queda condensada en esta igualdad: el marco
transportado por la conexi\'on de la totalidad conserva el dato situado en
el subespacio fijo.

Para una sub\`algebra de observables accesibles
\(\mathcal A_{\mathcal O}\), definimos
\begin{align}
 \mathcal R_{\rm vec}
 &=\{\psi:\mathscr H_\gamma\psi=\psi\},
 \label{eq:rm-obs-return-vector}\\
 \mathcal R_{\rm ray}
 &=\{\psi:\mathscr H_\gamma\psi=\ee^{i\theta}\psi
 \text{ para alg\'un }\theta\},
 \label{eq:rm-obs-return-ray}\\
 \mathcal R_{\rm obs}
 &=\{\psi:
 \langle\mathscr H_\gamma\psi,A\mathscr H_\gamma\psi\rangle
 =\langle\psi,A\psi\rangle
 \text{ para todo }A\in\mathcal A_{\mathcal O}\}.
 \label{eq:rm-obs-return-observable}
\end{align}

\begin{proposition}[Jerarqu\'ia vector--rayo--observable]
\label{prop:rm-obs-return-hierarchy}
Se cumple
\begin{equation}
 \boxed{
 \mathcal R_{\rm vec}
 \subseteq\mathcal R_{\rm ray}
 \subseteq\mathcal R_{\rm obs}.}
 \label{eq:rm-obs-return-hierarchy}
\end{equation}
\end{proposition}

\begin{proof}
Un vector fijo es un rayo fijo con fase cero. Si
\(\mathscr H_\gamma\psi=\ee^{i\theta}\psi\), las fases conjugadas se
cancelan en toda expectativa cuadr\'atica, por lo que todos los observables
accesibles conservan su valor.
\end{proof}

La monodrom\'ia non\'adica ocupa naturalmente esta jerarqu\'ia:
\begin{equation}
 g(k+9)=g(k),
 \qquad
 B_{k+9}\neq B_k\quad\hbox{en general}.
 \label{eq:rm-obs-phase-state-return}
\end{equation}
El retorno de fase pertenece a la lectura observable; el cambio de
\(B_k\) conserva el incremento de memoria. Una realizaci\'on inyectiva del
estado completo en la fibra permite elevar el criterio
\(\operatorname{Fix}(\mathscr H_\gamma)\) al retorno geneal\'ogico total.
En ausencia de esa inyectividad, el criterio afirma el retorno del dato
realizado.

\section{Banda de M\"obius y conjugaci\'on de orientaciones}

Sobre una palabra dodecaf\'asica
\(\mathbf t=(t_0,\ldots,t_{11})\), la involuci\'on geneal\'ogica es
\begin{equation}
 C_M(\mathbf t)=-\operatorname{rev}(\mathbf t),
 \qquad
 C_M^2=I.
 \label{eq:rm-obs-word-involution}
\end{equation}
Sea \(\ell_{\rm ret}\) un lazo de retorno enriquecido cuyo levantamiento
posee el car\'acter de signo
\begin{equation}
 \varepsilon_x:\langle\ell_{\rm ret}\rangle\simeq\mathbb Z
 \longrightarrow\{\pm1\},
 \qquad
 \varepsilon_x(\ell_{\rm ret})=-1.
 \label{eq:rm-obs-sign-character}
\end{equation}
El fibrado de intervalos asociado es
\begin{equation}
 \mathcal M_x
 =([0,1]\times[-1,1])/(0,u)\sim(1,-u).
 \label{eq:rm-obs-mobius-bundle}
\end{equation}

\begin{theorem}[Retorno relativo de M\"obius]
\label{thm:rm-obs-mobius-return}
Bajo el levantamiento compatible de \(C_M\), el lazo y el car\'acter de
\cref{eq:rm-obs-sign-character}, \(\mathcal M_x\) posee primera clase de
Stiefel--Whitney no nula. Una vuelta intercambia las orientaciones locales y
dos vueltas restauran la orientaci\'on de la secci\'on.
\end{theorem}

\begin{proof}
La funci\'on de transici\'on del cociente
\cref{eq:rm-obs-mobius-bundle} es la representaci\'on signo de
\(\pi_1(S^1)\). Su clase \(w_1\) es el generador no trivial de
\(H^1(S^1;\mathbb Z/2\mathbb Z)\). La primera iteraci\'on multiplica la
coordenada de fibra por \(-1\) y la segunda por \((-1)^2=1\).
\end{proof}

La formulaci\'on conserva tres retornos diferenciados: el cierre de fase
\(C_9\), el retorno de orientaci\'on en la doble cubierta y el incremento de
memoria en el levantamiento \(\mathbb Z\). Una vuelta puede restaurar el
rayo y cambiar el vector; dos vueltas restauran la orientaci\'on y pueden
conservar una profundidad distinta. La banda expresa precisamente la
diferencia entre publicaci\'on y genealog\'ia.

Los sectores par e impar de una funci\'on \(f\) sobre la cubierta doble son
\begin{equation}
 f^+=\frac12(f+C_M^*f),
 \qquad
 f^-=\frac12(f-C_M^*f).
 \label{eq:rm-obs-even-odd-decomposition}
\end{equation}
La parte par desciende a la base; la parte impar es una secci\'on de la
l\'inea de signo. La memoria cuadr\'atica reside en el sector par, mientras
orientaci\'on y carga lineal residen en el sector impar.

\section{Conjugaci\'on antiunitaria y orientaci\'on temporal}

La involuci\'on de palabra \(C_M\), la inversi\'on central del atlas y la
conjugaci\'on antiunitaria de la realizaci\'on f\'isica poseen dominios
distintos. Sea \(C\) la conjugaci\'on antiunitaria sobre una realizaci\'on de
Hilbert, sea \(H=H^*\) el Hamiltoniano y sea \(J_E\) la estructura compleja
en la realificaci\'on. Supongamos
\begin{equation}
 CJ_EC^{-1}=-J_E,
 \qquad
 CHC^{-1}=H,
 \qquad
 J_EH\subset HJ_E,
 \label{eq:rm-obs-antiunitary-hypotheses}
\end{equation}
con \(J_E^*=-J_E\), \(J_E^2=-I\) y preservaci\'on del dominio de \(H\).
La conmutaci\'on expresa la linealidad compleja de la din\'amica respecto de
la estructura \(J_E\).

\begin{proposition}[Reversi\'on antiunitaria de la evoluci\'on]
\label{prop:rm-obs-antiunitary-time}
Para
\begin{equation}
 U(t)=\exp(-J_EHt/\hbar_{\rm int}),
 \label{eq:rm-obs-unitary-evolution}
\end{equation}
se cumple
\begin{equation}
 \boxed{CU(t)C^{-1}=U(-t).}
 \label{eq:rm-obs-time-reversal}
\end{equation}
\end{proposition}

\begin{proof}
El generador \(-J_EH/\hbar_{\rm int}\) es antiadjunto. Las hip\'otesis de
\cref{eq:rm-obs-antiunitary-hypotheses} cambian su signo bajo conjugaci\'on.
El c\'alculo funcional y la unicidad del grupo unitario generado producen
\cref{eq:rm-obs-time-reversal}.
\end{proof}

Cuando, adem\'as, un operador de carga \(Q_{\rm ch}\) y un observable de
masa \(M\) satisfacen
\begin{equation}
 CQ_{\rm ch}C^{-1}=-Q_{\rm ch},
 \qquad
 CMC^{-1}=M,
 \label{eq:rm-obs-charge-mass-parity}
\end{equation}
las dos orientaciones de la banda se publican como portadores de carga
conjugada y memoria par compartida. La representaci\'on part\'icula--
antipart\'icula requiere el entrelazador que eleva \(C_M\) a \(C\) y
conserva la evoluci\'on. Las identidades
\cref{eq:rm-obs-word-involution,eq:rm-obs-time-reversal} proporcionan los
dos extremos tipados de ese entrelazador.

El lector temporal tr\'itico organiza los reg\'imenes por
\begin{equation}
 \begin{aligned}
 +1&\longmapsto\text{retorno, memoria y pasado},\\
 0&\longmapsto\text{umbral torsional y presente},\\
 -1&\longmapsto\text{apertura bidireccional y futuro}.
 \end{aligned}
 \label{eq:rm-obs-temporal-trit}
\end{equation}
La asignaci\'on es un lector de la estructura algebraica
\(J^2=-1,0,+1\). La promoci\'on a curvaturas de un espacio-tiempo exige una
conexi\'on y una carta m\'etrica; el lector temporal conserva mientras tanto
su significado exacto de retorno, umbral y apertura.

\section{Recuperabilidad y principio de Landauer}

El principio termodin\'amico de Landauer localiza el coste f\'isico del
borrado l\'ogicamente irreversible \cite{Landauer1961}.

Sea \((\Omega,2^\Omega,\mu)\) un espacio probabilizado finito, sea
\(X:\Omega\to\mathcal X\) la variable de estado completo y sea
\(q:\mathcal X\to Y\) un lector. La regla de la cadena proporciona
\begin{equation}
 \mathsf H(X)
 =\mathsf H(q(X))+\mathsf H(X\mid q(X)).
 \label{eq:rm-obs-entropy-chain}
\end{equation}
El segundo sumando es la informaci\'on contenida en la fibra del lector.

\begin{theorem}[Landauer nulo del transporte completo]
\label{thm:rm-obs-zero-landauer}
Si \(U:\mathcal X\to\mathcal X\) es una actualizaci\'on biyectiva, entonces
\begin{equation}
 \boxed{\mathsf H(X\mid U(X))=0.}
 \label{eq:rm-obs-zero-conditional-entropy}
\end{equation}
La informaci\'on l\'ogica borrada por la actualizaci\'on es nula. Para todo
protocolo f\'isico \(\mathcal P_U\) que la realice, la cota de Landauer y su
l\'imite reversible ideal satisfacen
\begin{equation}
 \boxed{
 Q_{\rm L}[\mathcal P_U]\geq0,
 \qquad
 \inf_{\mathcal P_U^{\rm rev}}Q_{\rm L}[\mathcal P_U]=0.}
 \label{eq:rm-obs-zero-landauer-cost}
\end{equation}
\end{theorem}

\begin{proof}
La inversa de \(U\) reconstruye \(X\) a partir de \(U(X)\); por tanto, la
distribuci\'on condicional est\'a concentrada en un solo estado. Su entrop\'ia
es cero y la operaci\'on l\'ogica conserva la informaci\'on. La cota
termodin\'amica resultante es \(Q_{\rm L}\geq0\); los protocolos
cuasiest\'aticos reversibles realizan el \'infimo ideal nulo. Esta afirmaci\'on
distingue la cota de la disipaci\'on concreta de un dispositivo.
\end{proof}

La salida proyectada \(q(U(X))\) conserva una parte de la informaci\'on y
deja en la fibra
\begin{equation}
 \mathsf H(X\mid q(U(X))).
 \label{eq:rm-obs-reader-forgotten-info}
\end{equation}
Sea \(M:\mathcal X\to\mathfrak M\) un registro de memoria y definamos el
lector enriquecido
\begin{equation}
 \widehat q(x)=(q(x),M(x)).
 \label{eq:rm-obs-enriched-reader}
\end{equation}

\begin{proposition}[Recuperabilidad mediante salida y memoria]
\label{prop:rm-obs-recoverability}
Si \(\widehat q\) es inyectivo sobre el soporte de \(X\), entonces
\begin{equation}
 \boxed{\mathsf H(X\mid q(X),M(X))=0.}
 \label{eq:rm-obs-recoverability}
\end{equation}
La cantidad \(\mathsf H(X\mid q(X))\) mide la memoria adicional requerida
para recuperar el estado desde la salida visible.
\end{proposition}

\begin{proof}
La inyectividad sobre el soporte proporciona una inversa determinista de
\(\widehat q(X)\) a \(X\). La entrop\'ia condicional es, por ello, cero. La
regla de la cadena localiza la diferencia entre el lector visible y el
lector enriquecido en la fibra de \(q\).
\end{proof}

Para el reinicio de un trit uniforme, la entrop\'ia l\'ogica es \(\log3\),
y la cota de Landauer toma la forma
\begin{equation}
 \boxed{Q_{\rm reset}\geq k_B\Theta\log3.}
 \label{eq:rm-obs-trit-reset}
\end{equation}
La identidad estructural del reloj,
\begin{equation}
 E_P=k_B\Theta_{\rm clk}\log3,
 \label{eq:rm-obs-planck-trit-landauer}
\end{equation}
eval\'ua esa cota en la temperatura cronol\'ogica declarada. El cuanto de
Planck aparece entonces como coste m\'inimo de reinicializar una decisi\'on
ternaria uniforme en esa carta t\'ermica.

El Landauer nulo y el coste \(k_B\Theta\log3\) describen operaciones
diferentes. El primero corresponde al transporte biyectivo del estado
enriquecido; el segundo corresponde a un reinicio que identifica tres
registros. Los costes de control, velocidad finita, preparaci\'on y
acoplamiento al entorno pertenecen a la realizaci\'on f\'isica del
dispositivo. La igualdad l\'ogica de coste m\'inimo cero conserva la
informaci\'on; una ecuaci\'on de movimiento para un rebote o un ciclo
cosmol\'ogico requiere, adem\'as, una din\'amica de frontera y sus balances.

\section{Medici\'on reversible, condicionamiento y reinicializaci\'on}

Las operaciones del proceso completo se ordenan como
\begin{equation}
 \boxed{\begin{gathered}
 \text{preparaci\'on}
 \longrightarrow
 \text{acoplamiento reversible}
 \longrightarrow
 \text{lectura del puntero}\\
 \longrightarrow
 \text{condicionamiento de la fibra}
 \longrightarrow
 \text{reinicializaci\'on}
 \end{gathered}}
 \label{eq:rm-obs-measurement-chain}
\end{equation}
Cada flecha posee un balance informacional propio:
\begin{enumerate}[label=\textup{(\roman*)}]
 \item La premedici\'on unitaria conserva pureza y norma en el sistema
 compuesto.
 \item La lectura publica \(q(X)\) y conserva el registro \(M(X)\) en el
 aparato.
 \item El condicionamiento selecciona la fibra relativa y normaliza su
 medida.
 \item El canal no selectivo suma todos los instrumentos y preserva la
 traza.
 \item La reinicializaci\'on identifica registros y activa la cota de
 Landauer correspondiente a la informaci\'on borrada.
\end{enumerate}

Esta cadena localiza el origen de la irreversibilidad operativa. El estado
global puede evolucionar isom\'etricamente mientras una lectura parcial
presenta entrop\'ia creciente. La informaci\'on accesible a una rama y la
informaci\'on conservada en las correlaciones globales son magnitudes
distintas. El observador acoplado forma parte de esas correlaciones y su
memoria constituye un componente f\'isico del estado conjunto.

\section{Confluencia entre observaci\'on, autoinercia y respuesta modular}

Los resultados anteriores se ensamblan en la cadena causal
\begin{equation}
 \begin{aligned}
 \text{estado APP--TRIT--TPK enriquecido}
 &\longrightarrow
 (\text{secci\'on},\text{observador},\text{memoria})\\
 &\longrightarrow
 (\text{fibra},\text{PVM},\text{instrumento})\\
 &\longrightarrow
 (\text{estado relativo},\text{holonom\'ia},\text{marco inducido})\\
 &\longrightarrow
 (\text{respuesta machiana},\text{Landauer tipado}).
 \end{aligned}
 \label{eq:rm-obs-causal-chain}
\end{equation}
La estructura modular del cap\'itulo anterior a\~nade el par
\begin{equation}
 (\widehat{\mathcal A},K_A),
 \label{eq:rm-obs-area-modular-pair}
\end{equation}
que conserva cantidad geom\'etrica y procedencia. El observador accede a una
compresi\'on de este par mediante su PVM; la memoria conserva el canal del
que procede el resultado; la autoinercia transporta el marco a lo largo de
la conexi\'on; la entrop\'ia relativa mide el coste de cambiar de estado
geneal\'ogico; y Landauer localiza el coste del borrado efectivo del
registro.

La relaci\'on con Everett se sit\'ua en el primer nivel de esta cadena: la
isometr\'ia conserva la superposici\'on de registros y el observador ocupa un
estado relativo. La relaci\'on con Mach se sit\'ua en el nivel de conexi\'on:
el marco y la respuesta local se inducen desde la totalidad compatible. La
banda de M\"obius se sit\'ua en el nivel de holonom\'ia: la orientaci\'on
puede cambiar mientras el rayo y los observables accesibles retornan. Las
tres lecturas comparten el estado enriquecido y conservan funciones
matem\'aticas diferentes.

\section{Dominio de validez y condiciones de realización física}

La fuerza probatoria se distribuye de la manera siguiente:
\begin{enumerate}[label=\textup{(\roman*)}]
 \item La correspondencia fibra--PVM--L\"uders, la representaci\'on de
 Kraus, el canal CPTP y la dilataci\'on son
 \status{resultados recuperados exactos} en los dominios declarados.
 \item La formulaci\'on del estado relativo y la conservaci\'on isom\'etrica
 del canal son \status{exactas}. La ontolog\'ia everettiana de ramas y su
 estabilidad macrosc\'opica pertenecen a una
 \status{interfaz interpretativa y din\'amica}.
 \item El criterio
 \(\operatorname{Fix}(\Hol)\), la jerarqu\'ia
 vector--rayo--observable y la factorizaci\'on machiana son
 \status{resultados recuperados exactos}. Su aplicaci\'on a una geometr\'ia
 gravitatoria exige el morfismo de realizaci\'on correspondiente.
 \item La banda de M\"obius es \status{exacta relativamente} al
 levantamiento, al lazo y al car\'acter de signo declarados. La
 interpretaci\'on part\'icula--antipart\'icula exige el entrelazador
 antiunitario de \cref{eq:rm-obs-antiunitary-hypotheses}.
 \item El Landauer nulo de una actualizaci\'on biyectiva, la recuperabilidad
 del lector enriquecido y la cota del reset ternario son
 \status{resultados recuperados exactos}. El selector cosmológico y la
 dinámica de hojas prolongan estos balances mediante las ecuaciones de
 evolución, las condiciones de empalme y la memoria torsional declaradas.
\end{enumerate}

\begin{proofstatusbox}
El cap\'itulo demuestra en sus dominios finitos la equivalencia de
condicionamientos, la conservaci\'on del canal global, los criterios de
retorno y de factorizaci\'on, y los balances l\'ogicos de Landauer. La
realización conjunta emplea el entrelazador isométrico entre la memoria TPK y
el aparato, la estabilidad dinámica del registro finito y el transporte de la
involución genealógica a la conjugación antiunitaria con sus observables pares
e impares.
\end{proofstatusbox}

\begin{falsifierbox}
Refutan los enunciados correspondientes: una fibra no medible; una familia
de proyectores que carezca de PVM com\'un; un condicionamiento de L\"uders
que difiera de \cref{eq:rm-obs-luders-fiber-equivalence}; operadores de
Kraus cuya suma de efectos difiera de la identidad; una dilataci\'on con
\(V_a^\dagger V_a\neq I\); un retorno vectorial exacto atribuido a un dato
fuera de \(\operatorname{Fix}(\Hol)\); una respuesta machiana variable dentro de una
fibra de \(b\); un lazo con car\'acter de signo trivial presentado como banda
de M\"obius; una conjugaci\'on que incumpla
\(CU(t)C^{-1}=U(-t)\); una actualizaci\'on biyectiva con entrop\'ia
condicional positiva; o un reset ternario uniforme cuyo balance omita la
informaci\'on \(\log3\). Tambi\'en refuta una promoci\'on cosmologica la
presentaci\'on del Landauer nulo como ecuaci\'on din\'amica de rebote.
\end{falsifierbox}
